\documentclass[12pt,a4paper]{article}
\usepackage{amsmath,amssymb,amsfonts}
\usepackage[margin=1in]{geometry}
\usepackage{enumitem}
\usepackage{fontspec}
\usepackage{polyglossia}
\setdefaultlanguage{bengali}
\setotherlanguage{english}
\newfontfamily\bengalifont[Script=Bengali]{Kalpurush}
\newfontfamily\englishfont[Script=Latin]{Times New Roman}
\newcommand{\mcq}[1]{{\bfseries উত্তর:} #1}
\newcommand{\exptext}[1]{{\bfseries ব্যাখ্যা:} #1}
\newcommand{\soln}[1]{{\bfseries সমাধান:} #1}
\title{\textbf{উচ্চতর গণিত ২য় পত্র — অধ্যায় ৭\\ বিপরীত ত্রিকোণমিতিক ফাংশন ও ত্রিকোণমিতিক সমীকরণ\\[4pt] \large সকল প্রশ্নের সমাধান ও ব্যাখ্যা (৯৪৩টি প্রশ্ন)}}
\author{}
\date{}

\begin{document}
\maketitle
\tableofcontents
\newpage

\input{batch01.tex}
\input{batch02.tex}
\input{batch03.tex}
\input{batch04.tex}
\input{batch05.tex}
\input{batch06.tex}
\input{batch07.tex}
\input{batch08.tex}
\input{batch09.tex}
\input{batch10.tex}
\input{batch11.tex}
\input{batch12.tex}
\input{batch13.tex}
\input{batch14.tex}
\input{batch15.tex}
\input{batch16.tex}
\input{batch17.tex}
\input{batch18.tex}
\input{batch19.tex}
\input{batch20.tex}
\input{batch21.tex}
\input{batch22.tex}
\input{batch23.tex}
\input{batch24.tex}
\input{batch25.tex}
\input{batch26.tex}
\input{batch27.tex}
\input{batch28.tex}
\input{batch29.tex}
\input{batch30.tex}
\input{batch31.tex}
\input{batch32.tex}

\end{document}
